%% Glossaire du jargon scientifique utilisé dans rapport_acc_nouveau_.tex
%% et sa version synthétique rapport_reduit_15p.tex
%% Clément Frassier — Projet HEART (UKM/UEL)

\documentclass[11pt]{article}
\usepackage[a4paper,margin=2.5cm]{geometry}
\usepackage[utf8]{inputenc}
\usepackage[T1]{fontenc}
\usepackage[french]{babel}
\usepackage{amsmath, amssymb}
\usepackage{enumitem}

\title{Glossaire du jargon scientifique\\\large Rapport « Apprentissage fédéré : détection du stress infirmier »}
\author{Clément Frassier}
\date{}

\begin{document}
\maketitle

\noindent Ce glossaire définit, dans l'ordre où ils apparaissent le plus naturellement
dans le rapport, tous les termes techniques employés -- métriques d'évaluation,
apprentissage fédéré, confidentialité différentielle, statistiques, et méthodologie du
jeu de données.

\section{Métriques d'évaluation}

\begin{description}[leftmargin=0em, style=nextline]

\item[Accuracy (précision brute)] Proportion de prédictions correctes sur l'ensemble
des exemples testés. Trompeuse sur un jeu de données déséquilibré : sur ce projet, une
prédiction ``stress'' systématique obtiendrait déjà $86\,\%$ d'accuracy sans avoir rien
appris, simplement parce que $86\,\%$ des fenêtres de test sont réellement du stress.

\item[Balanced-accuracy] Moyenne du rappel (voir ci-dessous) calculé séparément sur
chaque classe, plutôt que sur l'ensemble des exemples. Contrairement à l'accuracy, elle
n'est pas gonflée par une classe majoritaire : un modèle qui prédit toujours ``stress''
obtient $50\,\%$ de balanced-accuracy (niveau du hasard), quel que soit le déséquilibre
du jeu de données. C'est la métrique centrale de ce rapport.

\item[Rappel (\textit{recall})] Pour une classe donnée, proportion des exemples
réellement de cette classe que le modèle a correctement identifiés. Exemple : un rappel
``stress'' de $80\,\%$ signifie que le modèle détecte 8 épisodes de stress réels sur 10 --
la métrique cliniquement la plus parlante pour un système d'alerte (rater un épisode de
stress est plus grave que se tromper dans l'autre sens).

\item[Précision (\textit{precision}, au sens statistique)] À ne pas confondre avec
« précision brute » (accuracy) ci-dessus : ici, parmi toutes les fenêtres que le modèle a
prédites comme ``stress'', la proportion qui l'était réellement. Un modèle à haute
précision mais bas rappel prédit ``stress'' rarement, mais a raison quand il le fait.

\item[F1-score / F1 macro] Moyenne harmonique de la précision et du rappel, résumant les
deux en un seul chiffre. Le ``F1 macro'' moyenne ce score sur les deux classes (stress /
non-stress) à poids égal, comme la balanced-accuracy.

\item[AUC-ROC (\textit{Area Under the ROC Curve})] Mesure la capacité du modèle à
\emph{classer} les fenêtres de test par probabilité de stress croissante, indépendamment
d'un seuil de décision fixé. Vaut $0{,}5$ pour un modèle qui classe au hasard, $1{,}0$
pour un classement parfait. \textbf{Point important de ce rapport} : une AUC
\emph{significativement inférieure à $0{,}5$} ne signifie pas une absence de signal, mais
un classement \emph{partiellement inversé} -- le modèle a appris une relation
statistiquement réelle, mais dans le mauvais sens sur ces données précises.

\item[Seuil de décision (\textit{threshold})] Le modèle produit une probabilité continue
de stress (entre 0 et 1) ; le seuil est la valeur au-delà de laquelle on décide de
prédire ``stress''. Par défaut $0{,}5$. Le rapport teste aussi un balayage de ce seuil à
titre de \emph{diagnostic} (pour révéler un problème de calibration), en précisant
explicitement que ce n'est pas un protocole de déploiement validé.

\end{description}

\section{Apprentissage fédéré (FL)}

\begin{description}[leftmargin=0em, style=nextline]

\item[Apprentissage fédéré (\textit{Federated Learning}, FL)] Paradigme
d'entraînement collaboratif où plusieurs participants (``clients'') entraînent un modèle
commun sans jamais transmettre leurs données brutes à un serveur central -- seuls les
poids du modèle circulent.

\item[Client] Un participant à l'entraînement fédéré ; ici, une infirmière et son
bracelet connecté.

\item[Round] Un cycle complet d'échange : le serveur envoie le modèle courant aux
clients, chaque client l'entraîne localement, puis renvoie ses poids mis à jour au
serveur pour agrégation.

\item[Époque (\textit{epoch})] Un passage complet sur l'intégralité des données locales
d'un client pendant son entraînement, au sein d'un round.

\item[Non-IID (\textit{non independent and identically distributed})] Désigne des
données réparties de façon hétérogène entre les clients (ici, chaque infirmière a son
propre profil physiologique et sa propre composition de classes stress/non-stress) --
par opposition à des données ``IID'', où chaque client aurait un échantillon
statistiquement semblable des autres. Le Non-IID complique la convergence de
l'entraînement fédéré et est au cœur du message central de ce rapport.

\item[FedAvg] L'algorithme d'agrégation fédérée de base : le nouveau modèle global est la
moyenne des modèles des clients, pondérée par la quantité de données de chacun.

\item[FedProx] Variante de FedAvg qui ajoute un \textbf{terme proximal} à la fonction de
perte de chaque client, pénalisant l'éloignement entre les poids locaux et les poids
globaux -- corrige la ``dérive'' (\textit{client drift}) des clients en cas
d'hétérogénéité Non-IID, sans quoi l'entraînement peut diverger.

\item[Dérive des clients (\textit{client drift})] Phénomène où l'entraînement local de
chaque client, poursuivi sur ses propres données hétérogènes, écarte progressivement ses
poids de ceux des autres clients, rendant l'agrégation globale instable ou peu
performante.

\item[FedPer] Architecture de personnalisation où seule une partie du réseau
(l'``extracteur'', les premières couches) est partagée et agrégée entre clients ; la
dernière couche (la ``tête'' de classification) reste propre à chaque client et n'est
jamais transmise au serveur.

\item[SCAFFOLD] Mécanisme alternatif au terme proximal de FedProx pour corriger la
dérive Non-IID : chaque client maintient une ``variable de contrôle'' qui estime l'écart
entre la direction de mise à jour locale et globale, utilisée pour corriger le gradient
local à chaque étape.

\item[FedRep+SCAFFOLD] Dans ce rapport, une architecture combinant une tête de
classification à deux couches (au lieu d'une seule pour FedPer), une confidentialité
différentielle restreinte au seul extracteur partagé, et SCAFFOLD comme mécanisme de
correction de dérive -- explorée hors contrainte de taille TinyML, à titre de comparaison
haute plutôt que de cible de déploiement réel.

\item[Checkpoint] Une sauvegarde des poids d'un modèle à un instant donné de
l'entraînement (ici, typiquement à la fin des 30 rounds), permettant de le réévaluer ou
de le réutiliser sans tout ré-entraîner.

\item[Graine (\textit{seed})] Valeur initiale d'un générateur de nombres aléatoires.
Fixer la graine rend un résultat stochastique reproductible ; ce rapport répète
systématiquement ses expériences sur 5 graines indépendantes pour vérifier qu'une
conclusion ne dépend pas d'un tirage aléatoire particulier.

\end{description}

\section{Confidentialité différentielle}

\begin{description}[leftmargin=0em, style=nextline]

\item[Confidentialité différentielle (\textit{Differential Privacy}, DP)] Garantie
mathématique formelle qu'un mécanisme (ici, l'entraînement) ne révèle pas, avec une
probabilité significative, si les données d'un individu précis ont participé ou non à
l'entraînement -- une protection contre la ré-identification, plus forte qu'une simple
anonymisation.

\item[$(\varepsilon, \delta)$-DP] La formalisation quantitative de la confidentialité
différentielle. $\varepsilon$ (epsilon) est le \textbf{budget de confidentialité} :
plus il est petit, plus la garantie est forte (moins d'information sur un individu ne
``fuit''). $\delta$ (delta) est une probabilité résiduelle, très petite, que la garantie
échoue complètement -- généralement fixée bien en dessous de $1/n$ où $n$ est la taille
du jeu de données.

\item[DP-SGD] Version de la descente de gradient stochastique (l'algorithme
d'entraînement standard des réseaux de neurones) modifiée pour garantir le
$(\varepsilon,\delta)$-DP : à chaque étape, (1) le gradient de chaque exemple est
``clippé'' (borné en norme) puis (2) un bruit aléatoire gaussien calibré est ajouté à la
moyenne des gradients.

\item[Clipping de gradient] L'étape (1) ci-dessus : borner la contribution de chaque
exemple individuel à une norme maximale fixée ($C$), pour que le bruit ajouté ensuite
suffise à masquer la contribution de n'importe quel exemple.

\item[Multiplicateur de bruit ($\sigma$)] Paramètre contrôlant l'amplitude du bruit
gaussien ajouté par DP-SGD. Plus $\sigma$ est grand, plus la confidentialité est forte
(mais potentiellement plus le modèle perd en précision) -- le rapport en teste 5 valeurs
different pour mesurer ce compromis confidentialité/utilité.

\item[Opacus] Bibliothèque logicielle (PyTorch) qui implémente DP-SGD et calcule le
budget $\varepsilon$ consommé au fil de l'entraînement.

\item[\textit{PrivacyEngine}] L'objet central d'Opacus, créé une fois par client et
maintenu tout au long de l'entraînement fédéré pour que le budget $\varepsilon$
s'accumule correctement sur l'ensemble des rounds.

\item[Accountant (comptable de confidentialité)] Le composant qui calcule, à partir de
l'historique des étapes d'entraînement bruitées, le budget $\varepsilon$ effectivement
consommé. Ce rapport utilise l'accountant \textbf{PRV} (\textit{Privacy Random
Variable}) d'Opacus, successeur du \textbf{moments accountant} original.

\end{description}

\section{Modèle et déploiement embarqué}

\begin{description}[leftmargin=0em, style=nextline]

\item[MLP (\textit{Multi-Layer Perceptron})] Type de réseau de neurones le plus simple :
une succession de couches ``entièrement connectées'' (chaque neurone d'une couche est
relié à tous les neurones de la couche suivante), sans mécanisme récurrent ni
convolution. Le modèle ``MLP Tiny'' de ce rapport n'a que 1\,362 paramètres.

\item[TinyML] Discipline visant à faire tourner des modèles d'apprentissage automatique
directement sur des microcontrôleurs très contraints (mémoire de quelques dizaines à
centaines de kilo-octets, pas de GPU) -- ici, un bracelet connecté.

\item[FP32] Format de nombre à virgule flottante sur 32 bits, la précision standard
d'entraînement d'un réseau de neurones.

\item[INT8] Format entier sur 8 bits. La \textbf{quantification dynamique} post-entraînement
convertit les poids FP32 en INT8 pour réduire la taille du modèle ($\times3{,}8$
théorique ici) en vue du déploiement sur microcontrôleur, au prix d'une légère perte de
précision numérique (mesurée dans ce rapport et montrée quasi nulle en FL).

\item[Hyperparamètre] Un paramètre de configuration de l'entraînement fixé \emph{avant}
l'apprentissage (taux d'apprentissage, taille de batch, nombre de rounds, etc.), par
opposition aux poids du modèle qui sont appris.

\end{description}

\section{Tests et vocabulaire statistiques}

\begin{description}[leftmargin=0em, style=nextline]

\item[$p$-value] Probabilité d'observer un écart au moins aussi grand que celui mesuré
si, en réalité, il n'y avait aucune différence entre les deux conditions comparées. Une
$p$-value faible (conventionnellement $<0{,}05$) est interprétée comme un écart
``statistiquement significatif'', c'est-à-dire peu probable dû au seul hasard
d'échantillonnage.

\item[Test de Welch] Test statistique comparant les moyennes de deux groupes, adapté
même quand leurs variances diffèrent -- utilisé dans ce rapport pour comparer, par
exemple, le centralisé et le FL sur un petit nombre de graines.

\item[Test de Mann-Whitney (U)] Test statistique ``non paramétrique'' comparant deux
groupes sans supposer que les données suivent une loi normale -- plus robuste que le
test de Welch quand l'effectif est petit, utilisé ici en complément systématique.

\item[Test de Wilcoxon (signé, apparié)] Test non paramétrique pour des données
\emph{appariées} (ex. la même graine, avant/après une intervention). Avec seulement 5
paires dont les différences vont toutes dans le même sens, ce test plafonne
mécaniquement à une $p$-value de $0{,}0625$ -- un artefact de petit échantillon
documenté explicitement dans ce rapport, pas un signe de résultat douteux.

\item[Cohen's $d$] Mesure de la ``taille d'effet'' : l'écart entre deux moyennes,
exprimé en unités d'écart-type, indépendamment de la taille de l'échantillon.
Complète une $p$-value pour vérifier qu'un résultat significatif reflète un écart
réellement important, pas seulement un grand nombre de mesures.

\item[Correction de Bonferroni] Ajustement du seuil de significativité quand plusieurs
tests statistiques sont menés simultanément (ici, 4 comparaisons) : le seuil habituel
($0{,}05$) est divisé par le nombre de tests, pour éviter de déclarer ``significatif'' un
résultat qui ne l'est que par hasard du fait de tester plusieurs hypothèses à la fois.

\item[Écart-type] Mesure de la dispersion des valeurs autour de leur moyenne ; rapporté
systématiquement ici (``moyenne $\pm$ écart-type sur 5 graines'') pour indiquer la
variabilité d'un résultat d'une graine à l'autre.

\end{description}

\section{Méthodologie du jeu de données}

\begin{description}[leftmargin=0em, style=nextline]

\item[Fuite de données (\textit{data leakage})] Situation où des informations du
jeu de test s'infiltrent, directement ou indirectement, dans l'entraînement -- ce qui
gonfle artificiellement la performance mesurée sans refléter une vraie capacité de
généralisation. Ce rapport documente et corrige une fuite de ce type (fenêtres glissantes
qui se chevauchent entre train et test).

\item[Split chronologique] Découpage train/test qui respecte l'ordre temporel des
données (les premières mesures en entraînement, les dernières en test), par opposition à
un tirage aléatoire des exemples -- nécessaire ici pour éviter la fuite de données
ci-dessus.

\item[\textit{Gap} anti-fuite] Une petite portion de données délibérément exclue entre
la fin du train et le début du test, pour garantir qu'aucune fenêtre de l'un ne partage
de données brutes avec une fenêtre de l'autre.

\item[Standardisation] Transformation des données pour qu'elles aient une moyenne nulle
et un écart-type de 1 sur chaque caractéristique -- une étape de prétraitement standard.
Ce rapport distingue la standardisation \textbf{locale} (statistiques calculées par
client, utilisée en FL) de la standardisation \textbf{globale} (statistiques calculées
sur l'ensemble des clients regroupés, utilisée pour le centralisé).

\item[Test-set mono-classe] Un jeu de test qui ne contient qu'une seule des deux classes
(uniquement du stress, ou uniquement du non-stress). Découverte méthodologique centrale
de ce rapport : 11 des 15 infirmières sont dans ce cas, ce qui gonfle ou déforme
artificiellement toute métrique moyennée sur les 15 clients.

\item[Déséquilibre de classe (\textit{class imbalance})] Situation où une classe est
beaucoup plus représentée que l'autre dans un jeu de données (ici, $86\,\%$ de stress
contre $14\,\%$ de non-stress sur le test-set groupé) -- rend l'accuracy brute peu
informative et motive l'usage de la balanced-accuracy.

\item[Surapprentissage (\textit{overfitting})] Phénomène où un modèle apprend des
détails spécifiques à son jeu d'entraînement plutôt qu'un signal généralisable,
dégradant sa performance sur des données nouvelles.

\end{description}

\end{document}
